
\chapterstar{Planteamiento del Problema}

El cultivo de jitomate, pilar de la agroindustria mexicana, enfrenta una crisis fitosanitaria caracterizada por la alta incidencia de enfermedades virales, fúngicas y bacterianas que limitan severamente su potencial productivo. A pesar del uso de agricultura protegida, persisten patógenos de difícil control como el virus del rugoso del tomate (ToBRFV) y el tizón tardío (Rodríguez Alvarado et al., 2011).

Para abordar este problema, la identificación visual es la primera línea de defensa, permitiendo detectar síntomas y daños específicos:

\begin{itemize}
    \item \textbf{Tizón tardío (\textit{Phytophthora infestans}):} Se manifiesta con manchas cafés oscuras de aspecto aceitoso y, en condiciones de humedad, un crecimiento algodonoso blanco en el envés.
    \item \textbf{Mosca blanca (\textit{Bemisia tabaci}):} Provoca el amarillamiento de las hojas por la succión de savia y la excreción de melaza, que genera un hollín negro o fumagina.
    \item \textbf{Fusarium (\textit{F. oxysporum}):} Causa marchitez progresiva; el signo interno distintivo es un tejido vascular de color café rojizo al cortar el tallo.
    \item \textbf{Cenicilla o Mildiu:} Presenta un moho/polvillo blanco característico en el haz o envés, haciendo que las hojas se enrollen en forma de cuchara.
    \item \textbf{Alternaria (\textit{A. solani}):} Identificable por manchas negras o cafés con círculos concéntricos (anillos).
    \item \textbf{Plagas masticadoras y minadoras:} Los daños son evidentes a la vista por las ``minas'' o caminos blancos serpenteantes que deja la mosca minadora en el follaje, o las perforaciones circulares en el fruto causadas por el gusano del fruto.
    \item \textbf{Virus del Rugoso (ToBRFV):} Genera una rugosidad severa en la superficie del fruto y mosaicos de color verde-dorado en las hojas (FAO, 2013).
\end{itemize}

Este fenómeno es generalizado en México, afectando tanto cultivos a cielo abierto como bajo invernadero. Sin embargo, la problemática es crítica en los principales estados productores: Sinaloa (líder nacional), Michoacán, San Luis Potosí, Baja California, Sonora, Hidalgo y Puebla. En regiones como la zona centro de Michoacán, la severidad de estas enfermedades ha provocado que unidades de producción dejen de ser rentables.

La afectación es multidimensional:

\begin{enumerate}[topsep=2pt, itemsep=4pt, parsep=0pt]
    \item \textbf{Fisiológica:} Los patógenos alteran el funcionamiento interno de la planta, causando clorosis, necrosis, deformación de frutos y aborto floral.
    \item \textbf{Productiva:} Se genera una brecha entre el potencial de rendimiento (hasta 400 t/ha en invernaderos tecnificados) y la producción real obtenida.
    \item \textbf{Tecnológica:} Existe una dependencia de agroquímicos que, ante diagnósticos erróneos, resultan ineficaces e incrementan la resistencia de las plagas.
\end{enumerate}

Quiénes se ven afectados son:

\begin{itemize}
    \item \textbf{Productores agrícolas:} Quienes sufren la pérdida total o parcial de su inversión y, en muchos casos, el abandono de la actividad agrícola.
    \item \textbf{Economía Nacional:} Siendo México el principal exportador de jitomate fresco a nivel mundial, cualquier baja en la sanidad pone en riesgo divisas que superan los 6{,}000 millones de dólares anuales.
    \item \textbf{Consumidores finales:} Enfrentan inestabilidad en los precios y escasez del producto en el mercado nacional.
\end{itemize}

Actualmente, gran parte del monitoreo en cultivos de jitomate continúa realizándose mediante inspecciones visuales manuales y recorridos físicos dentro de los terrenos agrícolas, lo que dificulta la detección temprana de anomalías en grandes extensiones de cultivo. Debido a ello, muchas enfermedades son identificadas cuando el daño ya se encuentra avanzado, incrementando las pérdidas económicas y reduciendo la efectividad de las acciones correctivas. Ante esta problemática, surge la necesidad de implementar nuevas herramientas tecnológicas que permitan automatizar el monitoreo.

\thispagestyle{empty}
